\documentclass[../../../include/open-logic-section]{subfiles}

\begin{document}
\olfileid[hi]{sth}{card-arithmetic}{expotough}

\olsection[कुछ सरलीकरण]{कार्डिनल घातांक के कुछ सरलीकरण}

यद्यपि $\canonord$ को परिभाषित करना कुछ जटिल था, उसका परिणाम कार्डिनल योग
और गुणन संबंधी उपयोगी परिणाम \olref[simp]{cardplustimesmax} है। किंतु
परिमितेतर घातांक का इतना सीधा सरलीकरण नहीं किया जा सकता। इसका कारण समझाने
के लिए हम एक ऐसे परिणाम से आरंभ करते हैं जो परिमित स्थिति के परिचित प्रतिरूप
का विस्तार करता है (यद्यपि उसका प्रमाण अमूर्तन के उच्च स्तर पर है):

\begin{prop}\ollabel{simplecardexpo}
$\cardexpo{\cardfont{a}}{\cardfont{b} \cardplus \cardfont{c}} =
\cardexpo{\cardfont{a}}{\cardfont{b}} \cardtimes
\cardexpo{\cardfont{a}}{\cardfont{c}}$ और
$\cardexpo{(\cardexpo{\cardfont{a}}{\cardfont{b}})}{\cardfont{c}} =
\cardexpo{\cardfont{a}}{\cardfont{b} \cardtimes \cardfont{c}}$, किन्हीं
कार्डिनलों $\cardfont{a},\cardfont{b},\cardfont{c}$ के लिए।
\end{prop}

\begin{proof}
पहले दावे के लिए किसी फलन $f\colon
(\cardfont{b}\disjointsum\cardfont{c})\to\cardfont{a}$ पर विचार करें। अब
प्रत्येक $\beta\in\cardfont{b}$ के लिए
$f_\cardfont{b}(\beta)=f(\beta,0)$ और प्रत्येक
$\gamma\in\cardfont{c}$ के लिए
$f_\cardfont{c}(\gamma)=f(\gamma,1)$ परिभाषित करके ``इसे विभाजित करें''।
प्रतिचित्रण $f\mapsto(f_{\cardfont{b}}\times f_\cardfont{c})$ एक
!!a{bijection} है
$\funfromto{\cardfont{b} \disjointsum \cardfont{c}}{\cardfont{a}} \to
(\funfromto{\cardfont{b}}{\cardfont{a}} \times
\funfromto{\cardfont{c}}{\cardfont{a}})$. 

दूसरे दावे के लिए किसी फलन $f\colon\cardfont{c}\to
(\funfromto{\cardfont{b}}{\cardfont{a}})$ पर विचार करें; अतः प्रत्येक
$\gamma\in\cardfont{c}$ के लिए हमारे पास कोई फलन
$f(\gamma)\colon\cardfont{b}\to\cardfont{a}$ है। अब प्रत्येक
$\tuple{\beta,\gamma}\in\cardfont{b}\times\cardfont{c}$ के लिए
$f^*(\beta,\gamma)=(f(\gamma))(\beta)$ परिभाषित करें। प्रतिचित्रण
$f\mapsto f^*$ एक !!a{bijection} है
$\funfromto{\cardfont{c}}{(\funfromto{\cardfont{b}}{\cardfont{a}})}
\to \funfromto{\cardfont{b} \cardtimes \cardfont{c}}{\cardfont{a}}$. 
\end{proof}

अब अनंत कार्डिनलों से काम करते समय हम $\cardexpo{\cardfont{a}}{\cardfont{b}}$
का परिकलन करने का कोई सरल ढंग \emph{चाहेंगे}। इस दिशा में यह एक अच्छा कदम है:

\begin{prop}\ollabel{cardexpo2reduct}
यदि $2\leq\cardfont{a}\leq\cardfont{b}$ और $\cardfont{b}$ अनंत है, तो
$\cardexpo{\cardfont{a}}{\cardfont{b}}=\cardexpo{2}{\cardfont{b}}$।
\end{prop}

\begin{proof}
\begin{align*}
	\cardexpo{2}{\cardfont{b}} &\leq 
	\cardexpo{\cardfont{a}}{\cardfont{b}}\text{, क्योंकि $2 \leq \cardfont{a}$}\\
	&\leq \cardexpo{(2^\cardfont{a})}{\cardfont{b}}
	\text{, \olref[opps]{lem:SizePowerset2Exp} से}\\
	&= \cardexpo{2}{\cardfont{a} \cardtimes \cardfont{b}}
	\text{, \olref{simplecardexpo} से} \\
	&= \cardexpo{2}{\cardfont{b}}
	\text{, \olref[simp]{cardplustimesmax} से}
\end{align*}
\end{proof}

हमें वास्तव में इसके और सरलीकरण की अपेक्षा नहीं करनी चाहिए, क्योंकि
\olref[card-arithmetic][opps]{lem:SizePowerset2Exp} से
$\cardfont{b}<\cardexpo{2}{\cardfont{b}}$। किंतु यह हमें नहीं बताता कि
$\cardfont{b}<\cardfont{a}$ होने पर
$\cardexpo{\cardfont{a}}{\cardfont{b}}$ के बारे में क्या कहा जाए। निस्संदेह
यदि $\cardfont{b}$ \emph{परिमित} है, तो हम जानते हैं कि क्या करना है।

\begin{prop}
यदि $\cardfont{a}$ अनंत और $n\in\omega$, तो
$\cardexpo{\cardfont{a}}{n}=\cardfont{a}$।
\end{prop}

\begin{proof}
$\cardexpo{\cardfont{a}}{n}=\cardfont{a}\cardtimes\cardfont{a}
\cardtimes\ldots\cardtimes\cardfont{a}=\cardfont{a}$,
\olref[simp]{cardplustimesmax} से।
\end{proof}
\noindent 
इसके अतिरिक्त कुछ अन्य स्थितियों में हम
$\cardexpo{\cardfont{a}}{\cardfont{b}}$ के आकार को नियंत्रित कर सकते हैं:

\begin{prop}
यदि $2\leq\cardfont{b}<\cardfont{a}\leq
\cardexpo{2}{\cardfont{b}}$ और $\cardfont{b}$ अनंत है, तो
$\cardexpo{\cardfont{a}}{\cardfont{b}}=\cardexpo{2}{\cardfont{b}}$।
\end{prop}

\begin{proof}
$\cardexpo{2}{\cardfont{b}}\leq \cardexpo{\cardfont{a}}{\cardfont{b}}
\leq \cardexpo{(\cardexpo{2}{\cardfont{b}})}{\cardfont{b}} =
\cardexpo{2}{\cardfont{b}\cardtimes\cardfont{b}} =
\cardexpo{2}{\cardfont{b}}$, \olref{cardexpo2reduct} की ही रीति से।
\end{proof}
\noindent 
किंतु इस बिंदु के आगे विषय कहीं अधिक सूक्ष्म हो जाता है।

\end{document}
